% Template for abstracts submitted to ILAS 2020
% 1. Fill in the presenter's information (name, affiliation, email
% address) and talk information (title of presentation, abstract).
% 2. Compile to make sure there are no errors.
% 3. Save the .tex file for your abstract with a distinctive name,
% ideally "FamilynameGivenname.tex".
% 4. Upload your .tex file to https://clr.ie/129557
%       
% * Please keep your LaTeX commands simple, and avoid the use of
%    user-defined macros.
% * Include details of co-authors and funding acknowledgements after the abstract.
% * Upload only the LaTeX file (with .tex extension). 

\documentclass[a4paper]{amsart} 
\usepackage{amssymb} 
\usepackage{hyperref}
\pagestyle{empty}
\date{} 

%% IGNORE THE NEXT 4 LINES
\makeatletter % 
\def\labelenumi{\theenumi.}
\def\theenumi{\@arabic\c@enumi}
\makeatother
%% STOP IGNORING NOW

\begin{document}
	
	\title{Linearizations of rational matrices from general representations}
	\author{María C. Quintana} %Presenting author only
	\address{Aalto University, Finland.} % Affiliation of presenting author
	\email{maria.quintanaponce@aalto.fi} % Email address of presenting author only
	
	
	\maketitle
	
	% The abstract  ***no more than one page***
	
	Given a rational matrix $R(\lambda)$, the Rational Eigenvalue Problem (REP) consists of finding scalars $\lambda_0$ (eigenvalues) such that there exist nonzero constant vectors $x$ and $y$ (eigenvectors) satisfying 
	\[
	R(\lambda_0)x=0 \quad\mbox{and}\quad y^{T}R(\lambda_0)=0,
	\]
	under the regularity assumption $\text{det}\,R(\lambda)\not\equiv 0$. The numerical solution of REPs is recently getting a lot of attention from the numerical linear algebra community since REPs appear directly from applications or as approximations to arbitrary nonlinear eigenvalue problems. Rational matrices also appear in linear systems and control theory.
	
	Nowadays, a competitive method for solving REPs is linearization. Linearization transforms the REP into a generalized eigenvalue problem in such a way that the pole and zero information of the corresponding rational matrix is preserved. In this work, we construct a new family of linearizations of rational matrices $R(\lambda)$  written in the general form  $$R(\lambda)= D(\lambda)+C(\lambda)A(\lambda)^{-1}B(\lambda),$$ where $A(\lambda)$, $B(\lambda)$, $C(\lambda)$ and $D(\lambda)$ are polynomial matrices, with $A(\lambda)$ regular.
	Such representation always exists and are not unique. The new linearizations are constructed from linearizations of the polynomial matrices $D(\lambda)$ and $A(\lambda)$, where each of them can be represented in terms of any polynomial basis. In particular, the block minimal bases linearizations for polynomial matrices in \cite{BKL} will be our main tool for building linearizations of rational matrices in the sense of \cite{local}. In addition, we show how to recover eigenvectors, when $R(\lambda)$ is regular, and  minimal bases and minimal indices, when $R(\lambda)$ is singular, from those of their linearizations in this family. Finally, we show by example how the theory developed in this work can be used for solving (scalar) rational equations of the form
	$$
	\dfrac{c(\lambda)}{a(\lambda)}=\dfrac{d(\lambda)}{b(\lambda)},
	$$
	where $a(\lambda)$, $b(\lambda)$, $c(\lambda)$ and $d(\lambda)$ are nonzero scalar polynomials.
	

	
	%% Coauthor details here (optional - delete if not applicable)
	% and funding acknowledgment (optional - delete if not applicable)
	\emph{This is joint work with Javier Pérez (University of Montana, USA).\\
		Work (partially) supported by ``Ministerio de Econom\'ia, Industria y Competitividad (MINECO)'' of Spain and ``Fondo Europeo de Desarrollo Regional (FEDER)'' of EU through grants MTM2015-65798-P and MTM2017-90682-REDT, and the predoctoral contract BES-2016-076744 of MINECO.}
	
	
	%% IGNORE THE NEXT 4 LINES which change the enumerate style from (1) to [1]
	\makeatletter % 
	\def\labelenumi{[\theenumi]}
	\def\theenumi{\@arabic\c@enumi}
	\makeatother
	%% STOP IGNORING NOW
	
	
\begin{thebibliography}{1}

		\bibitem{local} F.~M.~Dopico, S.~Marcaida, M.~C.~Quintana, P.~Van Dooren,
		\newblock Local linearizations of rational matrices with application to rational approximations of nonlinear eigenvalue problems, \textit{Linear Algebra Appl.}, 604 (2020), 441--475.
		
		\bibitem{BKL} F.~M. Dopico, P.~W. Lawrence, J.~Pérez and P.~Van Dooren. Block Kronecker linearizations of matrix polynomials and their backward errors. \textit{Numer. Math.}, 140(2) (2018), 373--426.

\end{thebibliography}


	
	
\end{document}






